\chapter{Local determinant and Specht contractions}

Let $S=(K,\Omega,\mathcal A,\nu,\phi)$ be a proper adelic curve over a
countable field of characteristic zero. Write $K_\omega$ for the completion
of $K$ at $\omega$. At an Archimedean place its absolute value is the usual
real or complex absolute value. All vector spaces below have finite dimension.
All adelic bundles in this chapter are Hermitian.
The fibre of a bundle $\overline A$ is $A_\omega=K_\omega\otimes_K A$,
with its prescribed norm. Duals carry operator norms, subspaces induced norms,
and tensor products Hilbert norms at infinity and epsilon norms elsewhere.

For $a,k\geq1$ and $d=ak$, let $L_{a,d}$ be the rectangular integral Specht
lattice of Definition~\ref{def:constant-specht}. It sits in the free word
module on the words of length $d$ in $a$ letters. Each word occurring in it
contains each letter $k$ times. Its Archimedean metric is the restriction of
the metric making all words orthonormal; at other places its norm is the
maximum of its coordinates in an integral lattice basis. Put
\[
 S_{a,d}(\overline A)=(\det\overline A)^{\otimes k}
                   \otimes\overline L^S_{a,d},\qquad a=\dim_K A.
\]
The determinant norm at infinity is the Gram determinant norm. At a finite
place, for $\eta\neq0$ it is
\[
 \|\eta\|_{\det A}=\inf_{x_1\wedge\cdots\wedge x_a=\eta}
                         \prod_{i=1}^a\|x_i\|,
\]
and the norm of zero is zero. These are the determinant norms of
\cite[Section 2.3]{bost2026}. This chapter expands the local estimates in
\cite[Lemma 6.1 and the proof of Theorem 6.2]{bost2026}, including the
trivially valued case and the identifications on completed fibres.

\section{Ultrametric bases and determinant norms}

\begin{lemma}[Nearly orthogonal bases]
\label{lem:contraction-nearly-orthogonal-basis}
\source{bost-adelic-2026:page-0025}
\lean{BostConjecture.Adelic.Contraction.exists_basis_coord_bound}
\leanok
Let $F$ be a complete non-Archimedean valued field, possibly trivially valued,
and let $A$ be an ultrametric normed $F$-space of dimension $a$. For every
$0<\alpha<1$ there is a basis $b_1,\ldots,b_a$ whose coordinate functionals
satisfy $\alpha\|b_i\||b_i^*(x)|\leq\|x\|$ for every $i,x$.
Together with the ultrametric triangle inequality this says
\[
 \alpha\max_i |z_i|\|b_i\|
 \leq\Big\|\sum_i z_i b_i\Big\|
 \leq\max_i |z_i|\|b_i\|.
\]
\end{lemma}
\begin{proof}
Induct on $a$. The zero-dimensional assertion has empty maximum zero, and
the one-dimensional assertion follows from homogeneity. For $a>1$, choose
$\alpha<\beta<1$ and a hyperplane $U$. A finite-dimensional subspace of a
normed space over a complete valued field is closed, so for any $v\notin U$
the distance $\delta=\inf_{u\in U}\|v-u\|$ is positive. Replace $v$ by
$v-u_0$ with $\|v-u_0\|<\delta/\beta$. This does not change its coset or
$\delta$, and gives $\delta\geq\beta\|v\|$.
By induction choose an $(\alpha/\beta)$-orthogonal basis of $U$.
For $x=u+tv$, the quotient norm gives
$\|x\|\geq |t|\delta\geq\beta|t|\|v\|$. The ultrametric inequality
then gives $\|u\|\leq\max(\|x\|,|t|\|v\|)\leq\beta^{-1}\|x\|$.
Every weighted coordinate of $u$ is at most
$(\beta/\alpha)\|u\|\leq\alpha^{-1}\|x\|$, and the last weighted
coordinate $|t|\|v\|$ is also at most $\alpha^{-1}\|x\|$.
This proves the lower bound. The upper bound is the ultrametric triangle
inequality. The construction uses an infimum over vectors, not a scalar
of a prescribed absolute value, so it applies to a trivial valuation.
\end{proof}

The following weighted estimate expands the determinant step on page 25 of
this paper. Its arbitrary weights make explicit the calculation used there
with $w_i=\alpha\|b_i\|$.

\begin{lemma}[Weighted determinant estimate]
\label{lem:contraction-weighted-determinant}
\source{bost-adelic-2026:page-0025}
\lean{BostConjecture.Adelic.Contraction.prod_mul_norm_basis_det_le}
\leanok
Let $b$ be a basis of a normed space over a non-Archimedean valued field.
Write $b_i^*$ for its coordinate functionals. If $w_i>0$ and
$w_i|b_i^*(x)|\leq\|x\|$ for every $i,x$, then
\[
 \left(\prod_i w_i\right)
 \left|\det\bigl(b_i^*(x_j)\bigr)_{i,j}\right|
 \leq\prod_j\|x_j\|.
\]
\end{lemma}
\begin{proof}
Expand the determinant as a signed sum over permutations. The ultrametric
inequality bounds its absolute value by the largest absolute value of a
summand. For a permutation $\sigma$, multiplying the coordinate bounds for
$b_i^*(x_{\sigma(i)})$ gives the asserted bound on that summand, since
$\prod_i\|x_{\sigma(i)}\|=\prod_j\|x_j\|$.
The signs have absolute value one. There is no factorial factor.
\end{proof}

\begin{lemma}[Determinant lower bound]
\label{lem:contraction-completed-determinant}
\uses{lem:contraction-weighted-determinant}
\source{bost-adelic-2026:page-0025}
\lean{BostConjecture.Adelic.Contraction.determinant_norm_basis_ge}
\leanok
Fix a non-Archimedean place $\omega$ and a basis $b$ of $A_\omega$.
Suppose $\alpha>0$ and $\alpha\|b_i\||b_i^*(x)|\leq\|x\|$ for all
$i,x$. Set $\eta_b=b_1\wedge\cdots\wedge b_a$. Then
\[
 \|\eta_b\|_{\det A,\omega}\geq
                  \alpha^a\prod_i\|b_i\|_\omega.
\]
Here $\eta_b$ is viewed in $K_\omega\otimes_K\det A$ through the canonical
exterior-power base-change isomorphism.
\end{lemma}
\begin{proof}
For any $x_1\wedge\cdots\wedge x_a=\eta_b$, its determinant of
coordinates in $b$ is one. Apply the weighted estimate with
$w_i=\alpha\|b_i\|>0$. It gives
$\alpha^a\prod_i\|b_i\|\leq\prod_i\|x_i\|$.
Taking the infimum over these decompositions proves the result for the
specified determinant norm. The exterior base-change map sends a wedge of
completed vectors to that same wedge, so this computation is on the actual
completed determinant line.
\end{proof}

\section{Epsilon tensors and balanced words}

For a finite nonempty family of normed $F$-spaces $A_j$, define
\[
 \|z\|_\varepsilon=
 \sup_{\ell_j\neq0}
 \frac{|(\bigotimes_j\ell_j)(z)|}{\prod_j\|\ell_j\|_*}.
\]
We use the convention that the tensor of an empty family is $F$ with its
absolute value. The following lemmas justify the tensor operations used below
without requiring an orthogonal basis to exist exactly.

\begin{lemma}[Tensor maps on completed fibres]
\label{lem:contraction-tensor-maps}
\source{bost-adelic-2026:page-0027}
\lean{BostConjecture.Adelic.Contraction.orthogonalTensor_map_norm_le_of_finite_le_infinite_eq}
\leanok
Let $(\overline A_j)$ and $(\overline C_j)$ be finite families of Hermitian
bundles and $f_j:A_j\to C_j$ linear maps. Fix a place $\omega$.
Suppose $\|f_{j,\omega}(x)\|\leq\|x\|$ for every $j,x$, with equality
when $\omega$ is infinite. Then, for every completed tensor $z$,
\[
 \|(\bigotimes_j f_{j,\omega})(z)\|\leq\|z\|.
\]
\end{lemma}
\begin{proof}
At an infinite place, each factor map preserves the inner product by
polarization. The tensor map therefore preserves the product inner product
on pairs of pure tensors. These span, so sesquilinearity proves preservation
on all vectors, and hence the norm assertion over both $\mathbb R$ and $\mathbb C$.
At a finite place, a testing tuple of functionals $(\lambda_j)$ pulls back to
$(\lambda_j f_j)$, with $\|\lambda_j f_j\|_*\leq\|\lambda_j\|_*$.
The defining epsilon bound on $z$ therefore bounds the evaluation by
$\|z\|_\varepsilon\prod_j\|\lambda_j\|_*$.
A zero functional contributes zero. Taking the supremum proves the assertion.
\end{proof}

\begin{lemma}[Transposes of completed contractions]
\label{lem:contraction-dual-map}
\source{bost-adelic-2026:page-0027}
\lean{BostConjecture.Adelic.Contraction.dualMap_norm_le}
\leanok
Let $f:A\to C$ be a linear map between bundles with their prescribed norms.
At any place where $\|f_\omega(x)\|\leq\|x\|$ for all completed $x$,
its transpose satisfies $\|f_\omega^\vee(\lambda)\|_*\leq\|\lambda\|_*$.
\end{lemma}
\begin{proof}
For every $x$, one has
$|\lambda(f_\omega(x))|\leq\|\lambda\|_*\|f_\omega(x)\|
\leq\|\lambda\|_*\|x\|$. Taking the supremum over nonzero $x$ proves
the claim. Dual scalar extension identifies the completed transpose with
this precomposition: on scalar tensors both evaluate as
$(r\otimes\lambda)(s\otimes f(x))=rs\lambda(f(x))$, and bilinearity
extends the equality to arbitrary completed vectors.
\end{proof}

\begin{lemma}[Completed tensor duality]
\label{lem:contraction-tensor-duality}
\uses{lem:contraction-nearly-orthogonal-basis}
\source{bost-adelic-2026:page-0005,bost-adelic-2026:page-0027}
\lean{ArakelovGeometry.NormFamily.orthogonalTensor_dual_norm}
\leanok
For a finite family of Hermitian bundles, the canonical map
$\bigotimes_j A_{j,\omega}^\vee\to(\bigotimes_j A_{j,\omega})^\vee$
is an isometry for the prescribed completed tensor and dual norms.
\end{lemma}
\begin{proof}
At an infinite place, product orthonormal bases and their coordinate
functionals are orthonormal bases on both sides, so the map is isometric.
For a finite place, choose an $\alpha$-orthogonal basis $b_{j,i}$ of each
factor, where $0<\alpha<1$. Put $w_{\mathbf i}=\prod_j\|b_{j,i_j}\|$
and let $m$ be the number of factors.
For $z=\sum c_{\mathbf i}\bigotimes_j b_{j,i_j}$, testing with coordinate
functionals gives
$\|z\|_\varepsilon\geq\alpha^m\max_{\mathbf i}|c_{\mathbf i}|w_{\mathbf i}$.
The reverse bound without $\alpha^m$ follows by evaluating the expansion
against arbitrary testing functionals and applying the ultrametric inequality.
Write $u=\sum d_{\mathbf i}\bigotimes_j b_{j,i_j}^*$.
Evaluation on the $b_{j,i_j}$ gives
$\|u\|_\varepsilon\geq\max|d_{\mathbf i}|/w_{\mathbf i}$.
Consequently $|u(z)|\leq\alpha^{-m}\|u\|_\varepsilon\|z\|_\varepsilon$.
Letting $\alpha\uparrow1$ proves that the duality map contracts.
Conversely, if its image has operator norm $C$, evaluation on product basis
vectors gives $|d_{\mathbf i}|\leq Cw_{\mathbf i}$.
Since $\|b_{j,i_j}^*\|\leq(\alpha\|b_{j,i_j}\|)^{-1}$, the upper
coordinate estimate on $u$ gives $\|u\|_\varepsilon\leq\alpha^{-m}C$.
Let $\alpha\uparrow1$ again. The coordinate basis also proves that algebraic
tensor duality is bijective. The empty family gives scalar multiplication,
for which the assertion follows directly from the absolute value.
\end{proof}

\begin{lemma}[Tensor evaluation along an isometric inclusion]
\label{lem:contraction-tensor-evaluation}
\uses{lem:contraction-tensor-maps,lem:contraction-dual-map,lem:contraction-tensor-duality}
\source{bost-adelic-2026:page-0027}
\lean{BostConjecture.Adelic.Contraction.tensorPowerPairing_norm_le}
\leanok
Let $i:A\to W$ be linear between Hermitian bundles and fix a place at which
$i_\omega$ is isometric. For a finite index set $J$, the map
\[
 A_\omega^{\otimes J}\longrightarrow((W_\omega^\vee)^{\otimes J})^\vee,
 \qquad x\longmapsto[u\mapsto u(i_\omega^{\otimes J}x)]
\]
is contracting on every completed vector.
\end{lemma}
\begin{proof}
The tensor of $i_\omega$ contracts by the tensor-map lemma. Evaluation into
the bidual contracts because $|\lambda(x)|\leq\|\lambda\|_*\|x\|$.
The transpose of the tensor-duality isometry contracts by the transpose
lemma. Composing these three maps gives precisely the displayed evaluation
map, as is verified on pure tensors and extended linearly.
\end{proof}

\begin{lemma}[Completed canonical dual regrouping]
\label{lem:contraction-dual-regroup}
\uses{lem:contraction-nearly-orthogonal-basis,lem:contraction-tensor-maps,lem:contraction-tensor-duality}
\source{bost-adelic-2026:page-0027}
\lean{BostConjecture.Adelic.Contraction.dualRegroupEquiv_norm_le}
\leanok
For a finite family of Hermitian bundles $\overline A_j$ and a finite index
set $J$, the canonical map
\[
 \bigotimes_j(A_{j,\omega}^\vee)^{\otimes J}
 \longrightarrow ((\bigotimes_j A_{j,\omega})^\vee)^{\otimes J}
\]
is contracting for the actual nested tensor norms at every place.
\end{lemma}
\begin{proof}
First exchange the two tensor indices. At infinity this preserves the product
inner product on pairs of double pure tensors. Such tensors span, so the
map is isometric on arbitrary vectors by sesquilinearity.
At a finite place, the following epsilon universal bound proves contraction.
If $f:\bigotimes_{j=1}^m X_j\to Z$ is linear, $Z$ is ultrametric and
$\|f(\bigotimes_jx_j)\|\leq C\prod_j\|x_j\|$, expand in product
vectors from $\alpha$-orthogonal bases. The ultrametric triangle inequality
and coordinate-functional lower bound give
$\|f(z)\|\leq C\alpha^{-m}\|z\|_\varepsilon$.
Let $\alpha\uparrow1$ to obtain $\|f(z)\|\leq C\|z\|_\varepsilon$.
For nested tensors, apply this assertion to each row while holding all other
rows fixed, replacing the product of entry norms in that row by the row's
epsilon norm. Then apply it to the outer tensor. Thus a bound on double pure
tensors extends with the same constant to every vector of the nested tensor.
Regrouping sends a double pure tensor to its transpose; its target norm is
at most the product of all entry norms. Taking $C=1$ proves contraction.
Finally apply tensor duality to each column. These maps are isometries at all
places, so their tensor contracts by the tensor-map lemma. The composite is
the stated canonical dual regrouping. All these formulas commute with scalar
extension because they agree on scalar multiples of double pure tensors.
\end{proof}

\begin{lemma}[The balanced-word estimate]
\label{lem:contraction-balanced-words}
\uses{def:constant-specht}
\source{bost-adelic-2026:page-0026}
\lean{BostConjecture.Adelic.Contraction.injective_freeLatticeMap_le}
\leanok
Let $F$ be a complete field with an ultrametric absolute value.
Let $b_1,\ldots,b_a$ be a basis of a finite-dimensional normed $F$-space $A$,
and let $c_1,\ldots,c_f$
be any integral basis of $L_{a,ak}$. Write
\[
 x=\sum_h z_hc_h\in F\otimes_{\mathbb Z}L_{a,ak}.
\]
Let $b^{\otimes ak}x$ denote substitution of $b_i$ for the $i$-th letter.
Then
\[
 \|b^{\otimes ak}x\|_\varepsilon
 \leq\left(\prod_{i=1}^a\|b_i\|\right)^k\max_h|z_h|.
\]
\end{lemma}
\begin{proof}
Every $c_h$ has integral coefficients in the word basis, and every word in
its support uses each letter $k$ times. An integer has absolute value at most
one in a non-Archimedean field, including a trivially valued field.
Each substituted word has epsilon norm at most
$\prod_{q=1}^{ak}\|b_{i_q}\|=(\prod_i\|b_i\|)^k$, by testing with
tuples of dual functionals. The supremum of these ultrametric seminorms
satisfies the ultrametric triangle inequality. Applying it first to the word
expansion of $c_h$, and then to $\sum_hz_hc_h$, proves the bound.
\end{proof}

\section{The determinant embedding on completed fibres}

\begin{lemma}[Scalar extension of the determinant embedding]
\label{lem:contraction-determinant-base-change}
\uses{def:canonical-pairing,def:constant-specht}
\source{bost-adelic-2026:page-0025}
\lean{BostConjecture.Geometry.projectiveDeterminantMap_baseChange_of_basis,BostConjecture.Geometry.freeDeterminantMap_baseChange}
\leanok
The canonical determinant embedding
$\iota_A:(\det A)^{\otimes k}\otimes(K\otimes_{\mathbb Z}L_{a,ak})
\to A^{\otimes ak}$ commutes with field extension. In particular, under
the canonical exterior, lattice, and tensor base-change identifications,
its extension to $K_\omega$ is the map
\[
 \eta_b^{\otimes k}\otimes x\longmapsto b^{\otimes ak}x
\]
for every basis $b$ of $A_\omega$.
\end{lemma}
\begin{proof}
Start with a basis $e$ over $K$. The formula for $\iota_A$ on
$\eta_e^{\otimes k}\otimes c_h$ has integral word coefficients.
Extending scalars therefore gives the same substitution formula in the
extended basis $e$. These vectors form a basis of its source, proving
compatibility on the whole scalar extension.
For another basis $b=eg$, the determinant tensor transforms by
$\eta_b=\det(g)\eta_e$. On a rectangular Specht vector, substitution by
$g$ acts by $\det(g)^k$: a generating polytabloid is a permutation of a
product of $k$ alternating determinant tensors, each transforming by
$\det(g)$. Linear combinations give the same identity on the entire
lattice after scalar extension. Thus both sides of the defining rule
transform by $\det(g)^k$, proving the formula for every $b$.
\end{proof}

\begin{lemma}[Lattice estimate in determinant norm]
\label{lem:contraction-normalized-lattice}
\uses{lem:contraction-nearly-orthogonal-basis,lem:contraction-completed-determinant,lem:contraction-balanced-words,lem:contraction-determinant-base-change}
\source{bost-adelic-2026:page-0025,bost-adelic-2026:page-0026}
\lean{BostConjecture.Adelic.Contraction.injective_freeLatticeMap_completed_le}
\leanok
Fix a non-Archimedean place $\omega$ and a basis $b^0$ of $A_\omega$.
Choose an integral basis $(c_h)$ of $L_{a,ak}$. Write the coordinates of
$x\in K_\omega\otimes_{\mathbb Z}L_{a,ak}$ as $x=\sum_hz_hc_h$. Then
\[
 \|(b^0)^{\otimes ak}x\|_\varepsilon
 \leq\|\eta_{b^0}\|_{\det A,\omega}^{k}\max_h|z_h|_\omega.
\]
Both norms are the prescribed norms on the completed fibres, transported
by their canonical scalar-extension maps.
\end{lemma}
\begin{proof}
For $0<\alpha<1$, choose an $\alpha$-orthogonal basis $b$ and write
$\eta_{b^0}=\lambda\eta_b$. Basis independence of the determinant embedding
gives $(b^0)^{\otimes ak}x=\lambda^k b^{\otimes ak}x$.
The balanced-word estimate and the actual determinant lower bound imply
\[
 \alpha^{ak}\|(b^0)^{\otimes ak}x\|_\varepsilon
 \leq |\lambda|_\omega^k
       \left(\alpha^a\prod_i\|b_i\|_\omega\right)^k\max_h|z_h|_\omega
 \leq\|\eta_{b^0}\|_{\det A,\omega}^{k}\max_h|z_h|_\omega.
\]
The norms and coordinates involving $b^0$ are independent of $\alpha$.
Continuity of $\alpha\mapsto\alpha^{ak}$ at one therefore proves the
claim by letting $\alpha\uparrow1$. The existence and determinant estimates
used here also hold for trivial valuations.
\end{proof}

\begin{lemma}[Norm of a determinant--Specht generator]
\label{lem:contraction-determinant-specht-generator-norm}
\uses{def:constant-specht,lem:contraction-nearly-orthogonal-basis}
\source{bost-adelic-2026:page-0025,bost-adelic-2026:page-0026}
\lean{BostConjecture.Adelic.Contraction.determinantSpechtBundle_norm_tprod}
\leanok
At every place, for completed determinant vectors
$\eta_1,\ldots,\eta_k\in\det(A_\omega)$ and a completed Specht vector
$x\in K_\omega\otimes_{\mathbb Z}L_{a,ak}$, one has
\[
 \| (\eta_1\otimes\cdots\otimes\eta_k)\otimes x\|_{S_{a,ak}(A),\omega}
 =\left(\prod_{q=1}^k\|\eta_q\|_{\det A,\omega}\right)\|x\|_{L,\omega}.
\]
The left side uses the canonical exterior, tensor, and lattice scalar-extension
identifications with the original determinant--Specht bundle.
\end{lemma}
\begin{proof}
Hilbert tensor norms multiply on pure tensors by the product formula for
their inner products. Epsilon tensor norms also multiply on pure tensors:
the upper bound follows by testing with arbitrary dual functionals.
For the reverse bound, given a nonzero vector $y$ and $0<\alpha<1$, an
$\alpha$-orthogonal basis supplies a coordinate functional $\lambda$ with
$|\lambda(y)|\geq\alpha\|y\|\|\lambda\|_*$. Indeed choose a coordinate
maximizing $|c_i|\|b_i\|$, use
$\|y\|\leq\max_i|c_i|\|b_i\|$ and
$\|b_i^*\|\leq(\alpha\|b_i\|)^{-1}$.
Apply such functionals to each factor and let $\alpha\uparrow1$.
If a factor is zero both sides vanish; an empty determinant tensor has norm
one. Apply pure-tensor multiplicativity first to the determinant factors,
then to their tensor and $x$. The canonical scalar-extension maps send these
vectors to the same pure tensors, giving precisely the asserted norms.
\end{proof}

\begin{lemma}[Finite-place determinant embedding]
\label{lem:contraction-finite-specht}
\uses{lem:contraction-normalized-lattice,lem:contraction-determinant-specht-generator-norm,lem:contraction-determinant-base-change}
\source{bost-adelic-2026:page-0025,bost-adelic-2026:page-0026}
\lean{BostConjecture.Adelic.Contraction.projectiveDeterminantMap_norm_finite}
\leanok
Fix $\omega\in\Omega\setminus\Omega_\infty$.
The actual completed determinant embedding satisfies, for all
$z\in S_{a,ak}(A)_\omega$,
\[
 \|\iota_{A,\omega}(z)\|_\varepsilon\leq\|z\|_{S_{a,ak}(A),\omega}.
\]
\end{lemma}
\begin{proof}
Choose any basis $b$ of $A_\omega$.
Since $(\det A_\omega)^{\otimes k}$ is a line generated by
$\eta_b^{\otimes k}$, every $z$ is uniquely
$\eta_b^{\otimes k}\otimes x$. Its actual tensor norm is
$\|\eta_b\|^k\|x\|_L$, by the determinant--Specht generator norm formula.
Scalar-extension compatibility and the normalized lattice bound give
\[
 \|\iota_{A,\omega}(z)\|_\varepsilon
 \leq\|\eta_b\|^k\|x\|_L=\|z\|.
\]
The lattice norm equals the maximum of its integral-basis coordinates,
including at a trivially valued place. Thus these are exactly the prescribed
source and target norms.
\end{proof}

\begin{lemma}[Real and complex determinant embedding]
\label{lem:contraction-infinite-specht}
\uses{def:constant-specht,lem:contraction-determinant-base-change}
\source{bost-adelic-2026:page-0025}
\lean{BostConjecture.Adelic.Contraction.projectiveDeterminantMap_norm_infinite}
\leanok
At an Archimedean place $\omega$, the map $\iota_{A,\omega}$ on completed
fibres is an isometry for the prescribed determinant, Specht, and Hilbert
tensor norms.
\end{lemma}
\begin{proof}
Identify $K_\omega$ isometrically with $\mathbb R$ or $\mathbb C$, and
choose an orthonormal basis $b$ of the completed space with its given Hilbert
norm. Its Gram determinant is one, hence $\|\eta_b\|_{\det A}=1$.
The product vectors $b_{i_1}\otimes\cdots\otimes b_{i_{ak}}$ are an
orthonormal tensor basis. Consequently substitution by $b$ identifies the
word inner product with this Hilbert tensor inner product, and restricts to
the specified Specht Gram inner product. For every $x$ it follows that
\[
 \|\iota_{A,\omega}(\eta_b^{\otimes k}\otimes x)\|
 =\|b^{\otimes ak}x\|=\|x\|_L
 =\|\eta_b^{\otimes k}\otimes x\|.
\]
Every source vector has this form. Exterior and tensor base change send the
displayed generators to the same wedges and tensors, so these inner products
are precisely the prescribed completed norms. The computation uses real
symmetric Gram matrices over $\mathbb R$ and Hermitian Gram matrices over
$\mathbb C$ and therefore proves both cases. In particular, the
unnormalized alternating determinant tensor has norm $\sqrt{a!}$; that
factor remains in the word Gram norm and is not removed from the lattice.
\end{proof}

\section{The canonical pairing and its height}

Let $(\overline E_j)_{j\in I}$ be a finite nonempty family of positive-rank
Hermitian bundles, put $W=\bigotimes_j E_j$, and give a positive-dimensional
subspace $i:V\hookrightarrow W$ its induced norms. Write $a_j=\dim E_j$
and $s=\dim V$. Let $d=sl=a_jk_j>0$ for positive integers $l,k_j$, and set
\[
 B_d=S_{s,d}(\overline V),\qquad
 D_d=\bigotimes_j S_{a_j,d}(\overline E_j^\vee).
\]
These use exactly the determinant lines and constant Specht bundles already
specified. Let $j_d:D_d\to(W^\vee)^{\otimes d}$ be the tensor of the
determinant embeddings for the $E_j^\vee$, followed by canonical dual
regrouping. The canonical pairing of Definition~\ref{def:canonical-pairing} is
\[
 \Psi_d:B_d\longrightarrow D_d^\vee,\qquad
 \Psi_d(b)(u)=\langle j_d(u),i^{\otimes d}\iota_V(b)\rangle.
\]

\begin{lemma}[The completed pairing formula]
\label{lem:contraction-pairing-base-change}
\uses{def:canonical-pairing,lem:contraction-determinant-base-change}
\source{bost-adelic-2026:page-0026}
\lean{BostConjecture.Geometry.canonicalPairing_baseChange}
\leanok
For every place, the scalar extension $\Psi_{d,\omega}$ is the same
composite on completed fibres:
\[
 \Psi_{d,\omega}=j_{d,\omega}^{\vee}\circ
     \operatorname{ev}\circ i_\omega^{\otimes d}\circ\iota_{V,\omega}.
\]
Here $\operatorname{ev}:W_\omega^{\otimes d}\to
((W_\omega^\vee)^{\otimes d})^\vee$ is tensor evaluation, and the target
is identified with $(D_d^\vee)_\omega$ by dual base change.
\end{lemma}
\begin{proof}
The determinant embeddings commute with scalar extension by the preceding
lemma. On decomposable tensors, regrouping merely exchanges the two finite
indices, so regrouping also commutes with scalar extension.
Evaluation commutes with extension because
$(r\otimes\lambda)(t\otimes x)=rt\lambda(x)$ for $r,t\in K_\omega$.
These identities prove the formula on determinant generators and lattice
basis vectors on both sides of the pairing. Such generators span $B_d$ and
$D_d$ after extension; bilinearity therefore proves it for all completed
vectors. Finite-dimensional dual base change is an isomorphism.
\end{proof}

\begin{lemma}[Local canonical pairing contraction]
\label{lem:local-pairing-contraction}
\uses{def:canonical-pairing,def:constant-specht,lem:contraction-finite-specht,lem:contraction-infinite-specht,lem:contraction-tensor-maps,lem:contraction-dual-map,lem:contraction-tensor-evaluation,lem:contraction-dual-regroup,lem:contraction-pairing-base-change}
\source{bost-adelic-2026:page-0026,bost-adelic-2026:page-0027}
\lean{BostConjecture.Adelic.Contraction.canonicalPairing_norm_le}
\leanok
For every place $\omega$ and every $b\in(B_d)_\omega$,
\[
 \|\Psi_{d,\omega}(b)\|_{D_d^\vee,\omega}\leq\|b\|_{B_d,\omega}.
\]
Equivalently, $\|\Psi_{d,\omega}\|\leq1$. This includes both real and
complex places and every non-Archimedean place, even a trivially valued one.
No nonvanishing or geometric semistability assumption is needed.
\end{lemma}
\begin{proof}
The completed inclusion $i_\omega$ is isometric by the definition of the
induced norms. The two determinant-embedding lemmas prove that
$\iota_{V,\omega}$ and each determinant embedding defining $j_{d,\omega}$
are contracting and are isometric at infinite places. Tensor functoriality
and the contraction of dual regrouping
therefore make $j_{d,\omega}$ and $i_\omega^{\otimes d}$ contracting.
For arbitrary completed vectors $b,u$, the completed pairing formula and
the tensor evaluation estimate yield
\[
 \begin{aligned}
 |\Psi_{d,\omega}(b)(u)|
 &\leq\|j_{d,\omega}(u)\|\,
       \|i_\omega^{\otimes d}(\iota_{V,\omega}(b))\|\\
 &\leq\|u\|\|b\|.
 \end{aligned}
\]
Taking the operator dual norm in $u$ proves the asserted pointwise bound;
taking the operator norm in $b$ proves its equivalent form.
\end{proof}

\begin{corollary}[Nonpositive height of the canonical pairing]
\label{cor:canonical-pairing-height}
\uses{lem:local-pairing-contraction}
\source{bost-adelic-2026:page-0006,bost-adelic-2026:page-0027}
\lean{BostConjecture.Adelic.Contraction.canonicalPairing_mapHeight_le_zero}
\leanok
The canonical pairing has nonpositive lower-integral map height, with the
height of the zero map defined to be $-\infty$. In particular, if
$\Psi_d\neq0$, its ordinary map height is nonpositive:
\[
 h(\Psi_d)=\int_\Omega\log\|\Psi_{d,\omega}\|\,d\nu(\omega)\leq0.
\]
\end{corollary}
\begin{proof}
For the zero map the assertion follows from its height being $-\infty$.
Assume now that the pairing is nonzero. A nonzero map remains nonzero after
extension of fields, so each completed
operator norm is positive. The local contraction gives
$\log\|\Psi_{d,\omega}\|\leq0$ at every place. The logarithmic operator
norm of a nonzero map between adelic bundles is integrable; integrating the
pointwise inequality proves the claim. If map height is defined instead by
the extended lower integral, the same conclusion follows directly from
monotonicity, without an integrability hypothesis.
\end{proof}
